% Part: incompleteness
% Chapter: theories-computability
% Section: computably-axiomatizable

\documentclass[../../../include/open-logic-section]{subfiles}

\begin{document}

\olfileid{inc}{tcp}{cax}

\olsection{نظریه‌های \printtoken{S}{axiomatizable}}

نظریهٔ~$\Th{T}$ را \emph{!!{axiomatizable}} می‌نامیم اگر مجموعه‌ای
محاسبه‌پذیر از اصول موضوعه چون~$A$ داشته باشد. (اینکه~$A$ مجموعه‌ای
از اصول موضوعه برای~$\Th{T}$ است یعنی
$T = \Setabs{!A}{A \Proves !A}$.) هر اصل‌موضوعه‌بندی «معقولی» از
اعداد طبیعی این ویژگی را خواهد داشت. به‌ویژه، هر نظریه‌ای با مجموعهٔ
متناهی اصول موضوعه، !!{axiomatizable} است.

\begin{lem}
فرض کنید~$\Th{T}$، !!{axiomatizable} باشد. آنگاه~$\Th{T}$،
!!{computably
enumerable} است.
\end{lem}

\begin{proof}
فرض کنید~$A$ مجموعه‌ای محاسبه‌پذیر از اصول موضوعه برای~$\Th{T}$
باشد. برای نیمه‌تصمیم‌گیری دربارهٔ اینکه آیا~$!A \in T$، کافی است به دنبال
!!a{derivation} از~$!A$ بر پایهٔ اصول موضوعه بگردیم.

به بیانی اندکی متفاوت، $!A$ عضو~$\Th{T}$ است اگر و تنها اگر فهرستی
متناهی از اصول موضوعه~$!B_1$، \dots، $!B_k$ در~$A$ و
!!a{derivation} در منطق مرتبه‌اول از
$(!B_1 \land \dots \land !B_k) \lif !A$ وجود داشته باشد. اما از پیش
می‌دانیم هر مجموعه‌ای با تعریفی به صورت «\dots وجود دارد به‌گونه‌ای
که \dots» !!{c.e.} است، مشروط بر آنکه «\dots» دوم محاسبه‌پذیر باشد.
\end{proof}

\end{document}
